\documentclass[11pt]{article}
\usepackage[margin=0.9in]{geometry}
\usepackage{graphicx,amsmath,amssymb,booktabs,hyperref,xcolor,tabularx}
\graphicspath{{/tmp/dd_k3_optD_figs/}}
\hypersetup{colorlinks=true,linkcolor=blue,urlcolor=blue}
\setlength{\parskip}{6pt}
\title{Option \textit{d}: endogenous (equidistributed) p-grid for the K=3 lookup table}
\author{Fixed-Point Factory --- adaptive-grid follow-up}
\date{June 2026}

\begin{document}
\maketitle

\begin{abstract}
\noindent
Following the brainstorm, we test endogenous p-grids that equidistribute the
121 grid points by a monitor function $m(p)$ instead of the current
logit-uniform spacing. Four monitors compared at 16 saved $G{=}11$ R4
fixed points:
\begin{itemize}
  \item \textbf{uniform\_logit} (baseline): $m(p) = 1$, i.e.\ logit-uniform.
  \item \textbf{alpha\_a\_v}: $m(p) = a_v(p)$, the lookup density itself
        (averaged over $u_k$). Equal-mass bins.
  \item \textbf{beta\_dmu}: $m(p) = |\partial \mu_{\rm avg} / \partial p|$,
        the average $\mu$-slope (equal $\mu$-variation per bin).
  \item \textbf{gamma\_arc}: $m(p) = \sqrt{1 + |\partial \mu|^2}$, arc-length.
\end{itemize}
Each monitor's grid is built once from a reference run (strict-$h{=}0$ at
1001 fine pts), the lookup is recomputed on the adaptive grid, linearly
interpolated back to the reference grid, and the error vs the reference
measured.

\textbf{Result.} \emph{Alpha and beta} (effectively the same monitor: the
lookup density $\propto |\partial \mu /\partial p|$) cut the average max
interpolation error \textbf{from 0.312 to 0.234, a 25\% reduction} at the
same G\textsubscript{p}=121. Arc-length is similar (-23\%). Median error is
essentially zero for all monitors (the lookup is exact at the grid points
themselves; the residual is purely between-point interp).

\textbf{Where do the 121 points go?} The alpha monitor places \textbf{51 of
121 points} in the central region $|p - 0.5| < 0.05$ for $\tau{=}0.2$
(vs the baseline's 3), reflecting that the lookup is steepest at p=0.5
in the low-$\tau$ regime. At $\tau{=}1$ alpha keeps a few central points
and spreads the rest toward the tails where the lookup has features.

\textbf{Practical recommendation.} Adopt the alpha-monitor adaptive p-grid
as the production tabulation choice. Easy retrofit: compute the lookup once
on a reference grid, build the adaptive grid via cumulative-density
inversion, tabulate on the adaptive grid going forward. 25\% accuracy
improvement at zero extra runtime per Phi.
\end{abstract}

\section{Monitor comparison}

\begin{center}
\includegraphics[width=\linewidth]{fig1_monitor_compare.png}
\end{center}

\textbf{Left}: max interpolation error vs the reference (G\textsubscript{p}=1001 logit-uniform)
for each monitor across the 16 saved R4 FPs.
\textbf{Right}: median error. The bars for alpha and beta are visually
identical because both monitors are proportional to the lookup density
at the FP (the density IS the derivative of $\mu$ in the relevant sense).

\section{Average improvement}

\begin{center}
\includegraphics[width=\linewidth]{fig2_avg.png}
\end{center}

Average max interpolation error across all 16 cells:
\begin{center}
\begin{tabular}{lrr}
\toprule
monitor & avg max$|d|$ & improvement vs uniform\\
\midrule
uniform\_logit & 0.312 & --- \\
alpha\_a\_v & 0.234 & \textbf{-25\%} \\
beta\_dmu & 0.234 & \textbf{-25\%} \\
gamma\_arc & 0.239 & -23\% \\
\bottomrule
\end{tabular}
\end{center}

\section{Where the points go}

\begin{center}
\includegraphics[width=\linewidth]{fig3_density.png}
\end{center}

Average number of grid points the monitor places in the narrow band
$|p - 0.5| < 0.05$. The alpha monitor reconcentrates: from 3 baseline
points to $\sim 27$ adaptive points in the central region (averaged
over all cells, including high-$\tau$ where the central region matters
less).

\section{How to deploy}

\begin{enumerate}
  \item At the start of a sweep cell, solve to convergence on the current
        (logit-uniform) p-grid.
  \item At the FP, compute $a_v(p)$ on a dense reference (e.g.\ 1001 pts)
        via strict-$h{=}0$ or Option B+.
  \item Build the cumulative $A(p) = \int_0^p a_v(p') dp'$; pick new
        grid points $p_i = A^{-1}(i/G_p \cdot A(1))$.
  \item Re-tabulate the lookup on the new grid and re-feed Phase B
        clearing. Re-solve if you want a fully self-consistent FP.
\end{enumerate}

For repeated FP evaluations at nearby $(\gamma, \tau)$ (e.g.\ a sweep),
the adaptive grid can be reused across neighbors --- the monitor varies
slowly with parameters.

\section{Why median is zero}

The lookup is computed strict-$h{=}0$ \emph{exactly at the grid points}
$p_i$. The error comes only from linear interpolation between them. For
any monitor that places at least one grid point near each peak/feature of
$\mu(p, u_k)$, the median over the dense reference is effectively zero
(the median p-value is between two grid points, where linear interp is
sharp). The max error captures the worst-case between-grid-point feature.

\section{Reproducibility}
\texttt{projects/REZN/solver\_code/highprec\_dd\_qd/dd\_k3\_optD\_adaptive\_p.py}
(builder + tester) and \texttt{dd\_k3\_optD\_figs.py} (figures).
Results in
\texttt{projects/REZN/solved\_fixed\_points/dd\_k3\_overnight/optD/}.

\end{document}
